\section{The Offense at Nazareth}

The people who found Jesus of Nazareth hardest to believe in were the people who knew his street address. In the synagogue of his own town, the congregation heard him and asked the question that has never stopped being asked: ``Is not this the carpenter's son?'' They knew his mother's name. They could point to his relatives in the room. And the evangelist records the result without softening it: ``And they were scandalized in his regard.''\endnote{Matthew 13:55--57, quoted from the Douay--Rheims translation (Challoner revision), a public-domain English witness; wording checked at the online text of the Douay--Rheims Bible, 24 July 2026 and re-checked 25 July 2026. The same episode appears in Mark 6:1--6. Scripture quotations throughout are brief and taken from this witness unless another is identified.}

Mark's telling is blunter still. There the neighbors do not even grant the dignity of a patronymic: ``Is not this the carpenter, the son of Mary, the brother of James, and Joseph, and Jude, and Simon? are not also his sisters here with us?''---a complete inventory of local connections, offered as a refutation. And Mark reports the consequence in a sentence that has troubled copyists and commentators ever since: ``And he could not do any miracles there, only that he cured a few that were sick, laying his hands upon them.''\endnote{Mark 6:3--5, Douay--Rheims, checked 25 July 2026. The printed witness carries a note on ``he could not,'' glossing it as unwillingness rather than incapacity; the annotation belongs to that edition, not to this article, and the exegetical question is not adjudicated here. Nothing in the argument depends on how the clause is resolved.} Whatever else that verse means, it records a place where the too-familiar particular became an obstacle instead of a door---and it is the first datum this article has to account for.

The scandal was not that he claimed too much for God. It was that he located God too precisely. A carpenter's son, from this town, of that family---the offense lay in the definite articles. Nathanael, a few miles away and some months earlier, had compressed the whole difficulty into one question: ``Can any thing of good come from Nazareth?''\,\endnote{John 1:45--46, Douay--Rheims. Philip's answer---``Come and see''---is the Gospel's own procedure for the question and is deliberately not an argument.}

This reflex has a long intellectual pedigree, and it deserves to be stated with more respect than Nazareth gave it. In its serious form it runs roughly as follows. God, if the word means anything, means the absolute: unconditioned, unbounded, equally related to every time and place. But the Christian faith---and Israel's faith before it---presents a God who chooses one nation among the nations, becomes one man in one province in one decade, and now communicates himself through particular washings and particular suppers performed with water from this tap and bread from that oven. Surely, the objection says, the particular is beneath the universal. A truth of geometry is nobler than a fact of geography. A God worth the name would deal in humanity, not in Galileans; in goodness, not in loaves. To tie the absolute to the particular looks like a category mistake at best and provincial self-flattery at worst. Theologians have long called the difficulty the scandal of particularity, and Scripture itself uses the same vocabulary: the crucified Messiah is ``unto the Jews indeed a stumblingblock, and unto the Gentiles foolishness.''\,\endnote{1 Corinthians 1:22--25, Douay--Rheims. The Greek \emph{skandalon} underlying ``stumblingblock'' is the root of the theological phrase. ``Scandal of particularity'' itself is a modern term of art in theology and biblical studies, not a scriptural or magisterial formula; this article uses it as a received label for the objection, not as a doctrinal category.}

This article takes the objection seriously and argues that it is wrong---wrong not because the universal is unworthy of God, but because the objection misunderstands what universals are, what particulars are, and what kind of rescue human beings actually needed. The argument moves in an order that matters. It first lets the objectors speak at full strength and at length, because the objection is not a mood but a tradition with a documented history: it can be read in its own words across eighteen centuries, and reading it that way turns out to be part of answering it. It then examines the metaphysical premise hidden inside the complaint, because that premise can be tested by reason without appeal to revelation. Only afterward does it turn to what Israel, the Incarnation, and the sacraments claim, because those claims are not conclusions of philosophy and will not be presented as such. Philosophy can show that the objection fails; only the thing itself shows that the particular God acted. The distinction between those two kinds of showing is kept throughout, and the terminal appendices record, claim by claim, where each warrant lies.
